\chapter{The FX Options Market}\label{ch:m2:the-fx-options-market}

In most option markets a trader asks for the price of an option at a strike.
In foreign exchange the question is different. A broker's screen for
three-month dollar--yen shows three numbers: an at-the-money volatility, a
25-delta \omterm{def:m2:the-fx-options-market:rr}{risk reversal} and a 25-delta \omterm{def:m2:the-fx-options-market:rr}{butterfly}. There is no strike on it, and
the numbers do not mean the same thing in every pair: whether ``25 delta'' is a
spot or a \omterm{def:m2:the-fx-options-market:delta}{forward delta}, whether it counts the premium, which strike is ``at
the money'', all depend on conventions that differ between dollar--yen and
euro--dollar. Two traders who price the same option from the same screen with
different conventions get different strikes and different prices. FX options
more than doubled their trading between the central banks' surveys of 2022 and
2025, to 7\% of all FX turnover, and every one of those trades begins by
agreeing on these conventions. This chapter explains them, shows how a smile is
built from the three quotes, and follows the flows that \omterm{def:m2:the-fx-options-market:barrier}{barrier options} send
into the spot market.

\section{Delta and at-the-money conventions}

The model is the Black--Scholes model with two interest rates, the
Garman--Kohlhagen model: for a pair BASEQUOTE with spot $S$, the quote
currency's rate $r_d$ and the base currency's $r_f$, the forward is
$F = S e^{(r_d - r_f)T}$ and a call on the base currency with strike $K$ is worth
$e^{-r_d T}\bigl(F N(d_1) - K N(d_2)\bigr)$, with the usual $d_{1,2}$. What changes
from Book 1 is the delta, which FX quotes in four ways.

\begin{definition}[Spot, forward and premium-adjusted delta]\label{def:m2:the-fx-options-market:delta}
The \emph{spot delta}\index{spot delta} of an FX option is the amount of the base
currency, per unit of notional, to trade in the spot market to hedge it:
$\phi e^{-r_f T} N(\phi d_1)$, with $\phi = 1$ for a call and $-1$ for a put. The
\emph{forward delta}\index{forward delta} is the amount to trade in the forward
market, $\phi N(\phi d_1)$. The \emph{premium-adjusted
delta}\index{premium-adjusted delta} corrects either for a premium paid in the
base currency, which is itself a position in that currency:
$\phi e^{-r_f T} (K/F) N(\phi d_2)$ (spot) or $\phi (K/F) N(\phi d_2)$ (forward).
\end{definition}

The adjustment matters. A euro call whose delta is 60\% and whose premium,
73\,669 euros per million, is paid in euros is hedged by buying only 526\,331
euros: its \omterm{def:m2:the-fx-options-market:delta}{premium-adjusted delta} is 52.63\%. Which delta a pair uses is a
convention, documented in the literature on FX smile construction: pairs made of
rich-country currencies use \omterm{def:m2:the-fx-options-market:delta}{spot deltas} up to one year and \omterm{def:m2:the-fx-options-market:delta}{forward deltas} beyond;
pairs with an emerging-market currency use \omterm{def:m2:the-fx-options-market:delta}{forward deltas}; and the delta is
premium-adjusted when the premium is paid in the base currency, as in dollar--yen,
dollar--Swiss franc and euro--yen, but not in euro--dollar or sterling--dollar,
where it is paid in dollars, the quote currency.

\begin{definition}[Delta-neutral straddle]\label{def:m2:the-fx-options-market:dns}
The \emph{delta-neutral straddle}\index{delta-neutral straddle} is the call and the
put at the strike where their deltas add to zero; its strike is the default
``at the money'' for short-dated FX options: $K = F e^{\sigma^2 T/2}$ with regular
deltas and $K = F e^{-\sigma^2 T/2}$ with premium-adjusted ones.
\end{definition}

\begin{example}[Three-month strikes]\label{ex:m2:the-fx-options-market:strikes}
With EURUSD at 1.1464, dollar and euro rates of 3.68\% and 2.00\% (simple, for
three months) and a volatility of 7\%, the forward is 1.15119 and the
\omterm{def:m2:the-fx-options-market:dns}{delta-neutral straddle} strike 1.15190. For USDJPY at 156.87, with a yen rate of
0.977\% and a volatility of 9.5\%, the forward is 155.8196; the straddle strike is
155.9955 with regular deltas and 155.6439 with premium-adjusted ones, a gap of
35 pips of 0.01 yen between two definitions of the same ``at the money''.
\end{example}

\section{Risk reversals, butterflies and the smile}

\begin{definition}[Risk reversal, strangle, butterfly]\label{def:m2:the-fx-options-market:rr}
The 25-delta \emph{risk reversal}\index{risk reversal} is the difference between
the implied volatility of the 25-delta call and that of the 25-delta put,
$\sigma_{25C} - \sigma_{25P}$; as a trade, it is a long call and a short put (or
the reverse). The 25-delta \emph{strangle}\index{strangle} is a long 25-delta call
and a long 25-delta put; its volatility quoted over the at-the-money volatility is
the 25-delta \emph{butterfly}\index{butterfly}.
\end{definition}

The \omterm{def:m2:the-fx-options-market:rr}{risk reversal} measures the smile's slope: negative when the base currency's
puts are dearer than its calls, as when a fall of the dollar against the yen is
feared more than a rise. The \omterm{def:m2:the-fx-options-market:rr}{butterfly} measures its curvature: how much the wings
cost over the centre. A broker who says a \omterm{def:m2:the-fx-options-market:rr}{risk reversal} is ``1.2, puts over''
means that the put's volatility exceeds the call's by 1.2 points
(\cref{fig:m2:the-fx-options-market:quotes}). The simplest way
to turn the three quotes into two wing volatilities, often used as a first
approximation, is
\begin{equation}\label{eq:m2:the-fx-options-market:simplified}
\sigma_{25C} = \sigma_{\mathrm{ATM}} + \mathrm{BF} + \tfrac12\mathrm{RR}, \qquad
\sigma_{25P} = \sigma_{\mathrm{ATM}} + \mathrm{BF} - \tfrac12\mathrm{RR},
\end{equation}
which reproduces the \omterm{def:m2:the-fx-options-market:rr}{risk reversal} exactly. The market's quoted \omterm{def:m2:the-fx-options-market:rr}{strangle} is a
slightly different object from the average of the two wings, and building a smile
that reproduces it exactly needs a calibration; the approximation is the one we
use.

\begin{omfigure}
\begin{tikzpicture}
\begin{axis}[width=9cm,height=4.6cm,
  xlabel={delta},ylabel={implied volatility (\%)},
  tick label style={font=\small},label style={font=\small},
  xmin=-35,xmax=35,ymin=8.9,ymax=10.8,grid=major,grid style={gray!25},
  xtick={-25,0,25},xticklabels={25-delta put,ATM,25-delta call},
  table/col sep=comma]
\addplot[omThm,very thick,mark=*,mark size=2pt] table[x=x,y=vol]{figdata/markets-2/19-the-fx-options-market/quotes.csv};
\draw[omProp,thick,{Stealth}-{Stealth}] (axis cs:29,9.2) -- (axis cs:29,10.4);
\node[omProp,font=\scriptsize,anchor=east,align=right] at (axis cs:28.5,10.35) {RR\\$-1.2$};
\draw[omDef,thick,dashed] (axis cs:-25,9.8) -- (axis cs:25,9.8);
\node[omDef,font=\scriptsize,anchor=south] at (axis cs:6,9.83) {ATM + BF = 9.8};
\end{axis}
\end{tikzpicture}

\omcaption{The three quotes of \cref{ex:m2:the-fx-options-market:jpy} as three points
of the USDJPY smile: the put wing 10.4\%, at the money 9.5\%, the call wing 9.2\%.
The \omterm{def:m2:the-fx-options-market:rr}{risk reversal} is the difference between the wings, the \omterm{def:m2:the-fx-options-market:rr}{butterfly} the height
of their average over the centre. Illustrative.}\label{fig:m2:the-fx-options-market:quotes}
\end{omfigure}

\begin{proposition}[From quotes to a smile]\label{prop:m2:the-fx-options-market:smile}
Given $\sigma_{\mathrm{ATM}}$, RR and BF and a delta convention, the 25-delta
volatilities of \cref{eq:m2:the-fx-options-market:simplified} and the strikes
$K_{25C}$, $K_{25P}$ at which the options with those volatilities have deltas of
$0.25$ and $-0.25$ give, with the at-the-money strike, three points of the smile
in strike; any interpolation through them, here a quadratic in $\log K$, gives a
volatility at every strike.
\end{proposition}

\begin{proof}
For regular deltas the strike follows in closed form from
$N(\phi d_1) = \phi\Delta e^{r_f T}$ (spot) or $\phi\Delta$ (forward); for
\omterm{def:m2:the-fx-options-market:delta}{premium-adjusted deltas} it is found numerically, on the side of the strike range
where the delta is monotone. Three distinct points determine the quadratic.
\end{proof}

\begin{omfigure}
\begin{tikzpicture}
\begin{axis}[width=11cm,height=5cm,
  xlabel={USDJPY strike},ylabel={implied volatility (\%)},
  tick label style={font=\small},label style={font=\small},
  xmin=145,xmax=166,ymin=8.8,ymax=11.8,grid=major,grid style={gray!25},
  legend columns=2,legend style={font=\footnotesize,at={(0.5,-0.3)},anchor=north,draw=none,fill=none,/tikz/every even column/.append style={column sep=8pt}},
  table/col sep=comma]
\addplot[omDef,very thick] table[x=k_reg,y=v_reg]{figdata/markets-2/19-the-fx-options-market/smile_usdjpy.csv};
\addplot[omThm,very thick,dashed] table[x=k_pa,y=v_pa]{figdata/markets-2/19-the-fx-options-market/smile_usdjpy.csv};
\legend{regular spot deltas,premium-adjusted spot deltas}
\end{axis}
\end{tikzpicture}

\omcaption{A three-month USDJPY smile built from the same three quotes (ATM 9.5\%,
25-delta \omterm{def:m2:the-fx-options-market:rr}{risk reversal} $-1.2$, \omterm{def:m2:the-fx-options-market:rr}{butterfly} 0.3) with two delta conventions. The
premium-adjusted convention, the market's for this pair, places the wing and
centre strikes lower, so the same volatility belongs to a different strike.
Illustrative quotes; data: the chapter's tutorial.}\label{fig:m2:the-fx-options-market:smile}
\end{omfigure}

\begin{example}[Two smiles from one screen]\label{ex:m2:the-fx-options-market:jpy}
USDJPY three months: ATM 9.5\%, \omterm{def:m2:the-fx-options-market:rr}{risk reversal} $-1.2$ (dollar puts over), \omterm{def:m2:the-fx-options-market:rr}{butterfly}
0.3. By \cref{eq:m2:the-fx-options-market:simplified} the 25-delta dollar call is
at 9.2\% and the put at 10.4\%. With regular \omterm{def:m2:the-fx-options-market:delta}{spot deltas} their strikes are 160.85
and 150.71; with premium-adjusted \omterm{def:m2:the-fx-options-market:delta}{spot deltas}, the pair's convention, 160.68 and
150.52. A trader who priced the 150.71 put at 10.4\% while the market meant 150.52
would misprice it by the smile's slope over 19 pips of strike
(\cref{fig:m2:the-fx-options-market:smile}).
\end{example}

\section{Premium currency and premium-adjusted delta}

The premium currency changes both the hedge and the quote. A dollar-based bank
that sells a USDJPY call and receives its premium in dollars already holds some of
the dollars its hedge requires: it buys fewer. The market therefore quotes the
pair's deltas premium-adjusted, and the premium-adjusted call delta has a maximum:
deep in the money it rises with the strike, then falls again, so two strikes
can share a delta and the convention is to take the one above the maximum
(\cref{fig:m2:the-fx-options-market:deltas}). The build does so.

\begin{omfigure}
\begin{tikzpicture}
\begin{axis}[width=11cm,height=4.8cm,
  xlabel={USDJPY strike},ylabel={call spot delta},
  tick label style={font=\small},label style={font=\small},
  xmin=120,xmax=200,ymin=0,ymax=1.05,grid=major,grid style={gray!25},
  legend columns=2,legend style={font=\footnotesize,at={(0.5,-0.3)},anchor=north,draw=none,fill=none,/tikz/every even column/.append style={column sep=8pt}},
  table/col sep=comma]
\addplot[omDef,very thick] table[x=k,y=reg]{figdata/markets-2/19-the-fx-options-market/deltas.csv};
\addplot[omThm,very thick,dashed] table[x=k,y=pa]{figdata/markets-2/19-the-fx-options-market/deltas.csv};
\draw[gray,dashed] (axis cs:141,0) -- (axis cs:141,1.05);
\node[gray,font=\scriptsize,anchor=west,align=left] at (axis cs:166,0.62) {dashed: maximum\\of the premium-\\adjusted delta};
\legend{regular,premium-adjusted}
\end{axis}
\end{tikzpicture}

\omcaption{Three-month USDJPY call deltas by strike at 9.5\% volatility. The regular
delta falls steadily from one; the \omterm{def:m2:the-fx-options-market:delta}{premium-adjusted delta} is lower and, deep in the
money, rises with the strike before falling, so a delta below its maximum belongs
to two strikes. Illustrative; data: the chapter's
tutorial.}\label{fig:m2:the-fx-options-market:deltas}
\end{omfigure} The consequence for a trader is practical: before computing a
single number, check the pair's delta type, the at-the-money definition and the
premium currency of the quote.

\section{Broker markets and barrier-driven flows}

\begin{definition}[Barrier option, knock-out, knock-in]\label{def:m2:the-fx-options-market:barrier}
A \emph{barrier option}\index{barrier option} is an option whose existence
depends on whether spot touches a level, the barrier, before expiry. A
\emph{knock-out option}\index{knock-out option} ceases to exist when the barrier is
touched; a \emph{knock-in option}\index{knock-in option} comes into existence
only then.
\end{definition}

Between dealers, vanilla FX options are quoted as at-the-money straddles, \omterm{def:m2:the-fx-options-market:rr}{risk
reversals} and \omterm{def:m2:the-fx-options-market:rr}{strangles} rather than as single strikes, through brokers and on
electronic platforms. Barriers and other exotics are sold to companies and
investors and hedged by the dealers in vanillas and spot. Barriers appeal to
buyers because they are cheaper than vanillas, and they matter to the spot
market because their hedges change abruptly at the barrier.

\begin{omfigure}
\begin{tikzpicture}
\begin{axis}[name=left,width=6.4cm,height=4.8cm,
  xlabel={EURUSD spot},ylabel={value (pips of notional)},
  tick label style={font=\small},label style={font=\small},
  xmin=1.10,xmax=1.17,ymin=-5,ymax=320,grid=major,grid style={gray!25},
  xticklabel style={/pgf/number format/.cd,fixed,precision=2},
  table/col sep=comma]
\addplot[omThm,very thick] table[x=s,y=pips]{figdata/markets-2/19-the-fx-options-market/barrier.csv};
\draw[omProp,dashed] (axis cs:1.12,-5) -- (axis cs:1.12,320);
\node[omProp,font=\scriptsize,anchor=south west] at (axis cs:1.121,250) {barrier};
\end{axis}
\begin{axis}[at={(left.east)},anchor=west,xshift=1.4cm,width=6.4cm,height=4.8cm,
  xlabel={EURUSD spot},ylabel={delta},
  tick label style={font=\small},label style={font=\small},
  xmin=1.10,xmax=1.17,ymin=-0.05,ymax=0.9,grid=major,grid style={gray!25},
  xticklabel style={/pgf/number format/.cd,fixed,precision=2},
  table/col sep=comma]
\addplot[omThm,very thick,const plot] table[x=s,y=delta]{figdata/markets-2/19-the-fx-options-market/barrier.csv};
\draw[omProp,dashed] (axis cs:1.12,-0.05) -- (axis cs:1.12,0.9);
\end{axis}
\end{tikzpicture}

\omcaption{A three-month EURUSD down-and-out call, strike 1.1464, barrier 1.1200,
volatility 7\%: its value (left) and its delta (right) by spot. The value falls
to zero at the barrier, but the delta does not: it stays near 0.63 until the
touch and then drops to zero. Illustrative; data: the chapter's
tutorial.}\label{fig:m2:the-fx-options-market:barrier}
\end{omfigure}

\begin{proposition}[The hedge at the barrier]\label{prop:m2:the-fx-options-market:jump}
A dealer short a \omterm{def:m2:the-fx-options-market:barrier}{knock-out option} of notional $N$ and delta-hedged holds $\Delta N$
of the base currency. When spot touches the barrier the option and its delta
vanish, and the dealer must trade $\Delta(B^+)\,N$ at once, in the direction that
pushes spot further through the barrier: for a down-and-out call, selling at the
barrier.
\end{proposition}

\begin{proof}
Before the touch the hedge is $\Delta N$ with $\Delta \to \Delta(B^+) > 0$ as spot
approaches $B$ from above; after it, the option is gone and the hedge must be zero.
A down-and-out call's hedge is long the base currency, so unwinding it sells at
the barrier, where spot has just fallen.
\end{proof}

This is why barriers are defended and attacked. Holders of knock-outs, and dealers
long them, buy near the barrier to keep it from trading; dealers short them know
that their own hedge unwinds will add to a move through it, and those who know
where large barriers sit may try to push spot to them. Stop orders may sit at the
same levels, for the same reason: they are levels many people watch.

\begin{dated}{2026-09}{Size of the market}\label{dat:m2:the-fx-options-market:size}
In April 2025 FX options were 7\% of global FX turnover, more than double their
2022 turnover, according to the BIS Triennial Survey. The survey does not break
options into vanillas and barriers.
\end{dated}

\section{Tutorial: from quotes to a smile}\label{tut:m2:the-fx-options-market:smile}

\begin{tutorial}
\textbf{Goal.} Price FX options, compute their deltas in the four conventions,
recover the 25-delta and at-the-money strikes from broker quotes, build the smile,
and compute the hedge a knock-out releases at its barrier. \textbf{End state:}
\cref{fig:m2:the-fx-options-market:smile,fig:m2:the-fx-options-market:barrier},
\cref{ex:m2:the-fx-options-market:strikes,ex:m2:the-fx-options-market:jpy} and the
numbers of the weekend problem.
\begin{tutsteps}
\item \textbf{Deltas and strikes}: the four delta conventions and the strike with
a given delta.
\omcode{code/firm/fxsmile/firm_fxsmile.py}{32}{59}{Spot, forward and premium-adjusted deltas, and the strike with a given delta.}\label{lst:m2:the-fx-options-market:delta}
\item \textbf{At the money and the smile}: the \omterm{def:m2:the-fx-options-market:dns}{delta-neutral straddle}, the
simplified wing volatilities and a quadratic smile.
\omcode{code/firm/fxsmile/firm_fxsmile.py}{62}{81}{At-the-money strike, wing volatilities and a three-point smile.}\label{lst:m2:the-fx-options-market:smile}
\item \textbf{Run} \texttt{fxopt\_demo.smile} for each convention,
\texttt{fxopt\_demo.barrier()} and \texttt{fig\_fxopt.py}.
\end{tutsteps}
\textbf{What to change next.} Replace the simplified formula by a calibration that
matches the market \omterm{def:m2:the-fx-options-market:rr}{strangle} exactly, and compare the wing strikes; then value the
knock-out with the smile's volatility at the barrier instead of a flat one.
\end{tutorial}

\section{Build: the quote-to-smile converter}\label{bld:m2:the-fx-options-market:fxsmile}

\begin{build}
\bfield{Purpose} The miniature firm's FX option desk receives broker quotes in
ATM, RR and BF, and must price single strikes and barriers consistently with them
and hedge them with the right amount of spot.
\bfield{Interface} \texttt{gk(s, k, t, rd, rf, vol, phi)}; \texttt{delta(\ldots,
kind)} with kind in spot, forward, spot\_pa, forward\_pa;
\texttt{strike\_from\_delta}; \texttt{atm\_dns}; \texttt{smile\_vols(atm, rr, bf)};
\texttt{quadratic\_smile(points)}; \texttt{down\_and\_out\_call};
\texttt{barrier\_delta}.
\bfield{Rules} Continuously compounded rates; the premium-adjusted call strike
taken above the delta's maximum; the simplified formula for the wings, flagged as
an approximation.
\bfield{Acceptance tests} \texttt{code/firm/fxsmile/tests/}: put--call parity;
strike-to-delta round trips in all four conventions for calls and puts; the
published premium-adjustment example; delta-neutral strikes; the smile's
anchors; the barrier's value at and far from the barrier.
\bfield{Stretch} Calibration to the market \omterm{def:m2:the-fx-options-market:rr}{strangle}; smile interpolation in delta
(SABR or vanna--volga, One Quant Book 5); barrier pricing under the smile.
\end{build}

\begin{omsources}
\item D.~Reiswich and U.~Wystup, ``FX volatility smile construction'', Centre for
Practical Quantitative Finance, Working Paper 20, Frankfurt School of Finance and
Management, 2010 (conventions after I.~Clark, \emph{Foreign Exchange Option
Pricing}).
\item Bank for International Settlements, Triennial Central Bank Survey 2025.
\end{omsources}

\section{Exercises}

\begin{exercise}[$\star$]\label{exo:m2:the-fx-options-market:1}
The EURUSD 25-delta \omterm{def:m2:the-fx-options-market:rr}{risk reversal} is $-0.5$ and the \omterm{def:m2:the-fx-options-market:rr}{butterfly} 0.2, ATM 7\%. Give the
wing volatilities by the simplified formula, and say which option is dearer.
\end{exercise}

\begin{exercise}[$\star$]\label{exo:m2:the-fx-options-market:2}
Which delta convention does the market use for EURUSD, USDJPY and USDBRL, and why
does the premium currency matter?
\end{exercise}

\begin{exercise}[$\star$]\label{exo:m2:the-fx-options-market:3}
A euro call has a \omterm{def:m2:the-fx-options-market:delta}{spot delta} of 60\% and costs 73\,669 euros per EUR~1 million,
paid in euros. How many euros does the seller buy to hedge?
\end{exercise}

\begin{exercise}[$\star\star$]\label{exo:m2:the-fx-options-market:4}
Give the three-month EURUSD 25-delta call and put strikes with regular spot and
\omterm{def:m2:the-fx-options-market:delta}{forward deltas}, using the quotes of \cref{exo:m2:the-fx-options-market:1}.
\end{exercise}

\begin{exercise}[$\star\star$]\label{exo:m2:the-fx-options-market:5}
Why is the premium-adjusted delta-neutral strike below the forward, and the
regular one above?
\end{exercise}

\begin{exercise}[$\star\star$]\label{exo:m2:the-fx-options-market:6}
Why is a down-and-out call cheaper than the vanilla call, and when is the
difference large?
\end{exercise}

\begin{exercise}[$\star\star\star$]\label{exo:m2:the-fx-options-market:7}
\emph{Coding.} With \texttt{strike\_from\_delta}, give the USDJPY 25-delta strikes
with \omterm{def:m2:the-fx-options-market:delta}{forward deltas}, and compare them with the spot and premium-adjusted ones of
\cref{ex:m2:the-fx-options-market:jpy}.
\end{exercise}

\begin{exercise}[$\star\star\star$]\label{exo:m2:the-fx-options-market:8}
\emph{Find the flaw.} ``The \omterm{def:m2:the-fx-options-market:rr}{risk reversal} is negative, so the market expects the
dollar to fall against the yen.'' Correct it.
\end{exercise}

\section{Problem: The Barrier Defence}

\begin{problem}[{Weekend problem --- what a knock-out sends into the spot market}]\label{pb:m2:the-fx-options-market:1}
A company buys from a dealer a three-month EURUSD down-and-out call on EUR~500
million, strike 1.1464, barrier 1.1200, spot 1.1464, with the rates of
\cref{ex:m2:the-fx-options-market:strikes} and a volatility of 7\%. The dealer
delta-hedges.

\medskip\noindent\textbf{Part I --- The option.}
\begin{enumerate}
\item Give the vanilla call's value, in pips of notional.
\item Give the knock-out's value, and the saving for the company.
\item Give the knock-out's delta now, and the dealer's hedge.
\item Why is the knock-out's delta larger than the vanilla's?
\item Give the delta just above the barrier, and the vanilla's there.
\end{enumerate}

\medskip\noindent\textbf{Part II --- At the barrier.}
\begin{enumerate}[resume]
\item What does the dealer hold just before spot touches 1.1200?
\item What must it trade at the touch, and in which direction?
\item How does that trade affect spot?
\item What would a dealer long such a barrier do as spot approaches it?
\item Why might others trade towards a known large barrier?
\end{enumerate}

\medskip\noindent\textbf{Part III --- Managing it.}
\begin{enumerate}[resume]
\item How can the dealer spread the unwind instead of doing it at the touch?
\item What does that cost, and what risk does it leave?
\item Why does the dealer's gamma near the barrier make daily hedging expensive?
\item How would a volatility smile change the knock-out's value?
\item What should the company understand about the option it bought?
\end{enumerate}

\medskip\noindent\textbf{Part IV --- Judgement.}
\begin{enumerate}[resume]
\item Is defending a barrier market manipulation?
\item Why do barriers cluster at round numbers?
\item What information about barriers should a dealer keep confidential?
\item State the \emph{named result}: the spot the dealer must sell at the barrier.
\item In one sentence: why do \omterm{def:m2:the-fx-options-market:barrier}{barrier options} move spot?
\end{enumerate}
\end{problem}

\section{Interview questions}

\begin{interviewq}[$\star$ \iqroles{trader, researcher}]\label{iq:m2:the-fx-options-market:1}
What do the at-the-money volatility, the 25-delta \omterm{def:m2:the-fx-options-market:rr}{risk reversal} and the \omterm{def:m2:the-fx-options-market:rr}{butterfly}
tell you about a currency pair's smile?
\end{interviewq}

\begin{interviewq}[$\star$ \iqroles{developer, trader}]\label{iq:m2:the-fx-options-market:2}
What is a \omterm{def:m2:the-fx-options-market:delta}{premium-adjusted delta}, and when is it used?
\end{interviewq}

\begin{interviewq}[$\star\star$ \iqroles{researcher}]\label{iq:m2:the-fx-options-market:3}
Given ATM, RR and BF quotes, how do you find the strike of the 25-delta put?
\end{interviewq}

\begin{interviewq}[$\star\star$ \iqroles{trader}]\label{iq:m2:the-fx-options-market:4}
You are short a large down-and-out call and spot is falling towards the barrier.
What do you do?
\end{interviewq}

\begin{interviewq}[$\star\star$ \iqroles{researcher, trader}]\label{iq:m2:the-fx-options-market:5}
Why is at-the-money in FX defined as the \omterm{def:m2:the-fx-options-market:dns}{delta-neutral straddle} rather than the
forward?
\end{interviewq}

\begin{interviewq}[$\star\star\star$ \iqroles{developer}]\label{iq:m2:the-fx-options-market:6}
Design a volatility-surface service for fifty currency pairs that accepts broker
quotes and serves volatilities by strike and expiry to pricing engines.
\end{interviewq}
